\year=2026
\month=10

\day=05

\setcounter{chapter}{1}
\chapter{Static Equilibrium \& Elasticity}
\section{Conditions \& Examples for Static Equilibrium}
\subsection{Rigid Body}

The rigid body depends on the \emph{physical properties} of the material

\subsection{Free-body diagrams}
\subsection{Equilibrium}

\textbf{Equilibrium:} the state of balance in a system
There is \emph{no linear motion} \& \emph{no rotational motion}

A body in equilibrium can be moving
but if so its linear \& angular velocities must be constant

The conditions of static equilibrium are
\begin{enumerate}
	\item The result force on the system \( F_{\text{net}} = 0 \)
	\item The result torque on the system \( \sum\tau = 0 \)
\end{enumerate}

\subsection{The first condition of Equilibrium}
The sum of the forces acting on the rigid body is zero,
\[
	\vec{F} = \vec{F_1} + \vec{F_2} \ldots = 0
\]
\begin{itemize}
	\item	the horizontal forces, \( \sum F_x = F_{\text{net}x} = 0 \)
	\item	the vertical forces,   \( \sum F_y = F_{\text{net}y} = 0 \)
\end{itemize}

\begin{figure}[htbp]
	\centering
	\begin{tikzpicture}[scale=4]
		\coordinate (O) at (0,0);

		\draw[thin, ->] (O) ++(-1,0) -- (O) -- ++(1,0) node[right] {$x$};
		\draw[thin, ->] (O) ++(0,-1) -- (O) -- ++(0,1) node[above] {$y$};

		\tikzset{
			vecstyle/.style={draw=green, thick, ->}
		}
		\draw[vecstyle] (O) -- ++(25:1) node[above] {$F_1$};
		\draw[vecstyle] (O) -- ++(143:1) node[above] {$F_3 = 29.0 N$};
		\draw[vecstyle] (O) -- ++(-90:1) node[right] {$F_2 = 22.0 N$};
	\end{tikzpicture}
	\label{fig: three forces}
\end{figure}
In figure,
the net force is
(if we take the direction to the right to be positive)
\[
	\sum F = 12 + 6.0 - 8.0 = 10 N
\]

\begin{align*}
	F_{2x} &= 0 \\
	F_{2y} &= -22.0 N \\
	F_{3x} &= -29.0\cos{37\deg} = -23.16 N \\
	F_{3y} &= 29.0\sin{37\deg} = 17.45 N \\
	\text{In equilibrium,}
	\sum F_{x} = 0 &\mathbin{\&}
	\sum F_{y} = 0 \\
	F_{1x} + F_{2x} + F_{3x} &= 0 \\
	F_{1x} + 0 - 23.16  &= 0 \\
	F_{1x} &= 23.16 N \\
	%
	F_{1y} + F_{2y} + F_{3y} &= 0 \\
	F_{1y} - 22.0 + 17.45 &= 0 \\
	F_{1y} &= 4.55 N \\
	%
	\text{Therefore,}\enspace
	F_{1} &= \sqrt{ {F_{1x}}^2 + {F_{1y}}^2 } \\
	F_{1} &= \sqrt{ {23.16}^2 + {4.55}^2 } = 23.6 N \\
	\text{The angle is found from,}\enspace
	\tan\theta &= \frac{F_{1y}}{F_{1x}},
	\theta = \tan^{-1} \frac{4.55}{23.16} = 11.115\deg
\end{align*}

Torque is a kind of "rotational force":
\( \vec\tau = \vec{r} \times \vec{F} \)

\textbf{Unit of torque}
\( \lvert \tau \rvert = r F_{\perp} = Nm \)

The torque is
\begin{itemize}
	\item \textbf{positive} when \textbf{counter-clockwise} rotation
	\item \textbf{negative} when \textbf{clockwise} rotation
\end{itemize}
